Deoarece automatul a fost proiectat având ieșiri de tip Moore (asociate stărilor, nu tranzițiilor), funcțiile de ieșire depind exclusiv de starea curentă ($Q_3 Q_2 Q_1 Q_0$).

Conform organigramei și tabelelor de stare, avem următoarele asocieri:
\begin{itemize}
    \item \textbf{ERR\_LED} este activ doar în starea \textbf{S7 (ERR)} - cod \texttt{0111}.
    \item \textbf{SPK\_EN} și \textbf{RX\_LED} sunt active doar în starea \textbf{S11 (LSTN)} - cod \texttt{1011}.
    \item \textbf{MIC\_EN} și \textbf{TX\_LED} sunt active doar în starea \textbf{S13 (TRSM)} - cod \texttt{1101}.
\end{itemize}

Mai jos sunt prezentate hărțile Karnaugh și ecuațiile rezultate pentru fiecare grup de ieșiri.

\subsection{Implementarea ieșirii de eroare}

Această ieșire semnalizează eșecul împerecherii. Este activă când $Q_3 Q_2 = 01$ și $Q_1 Q_0 = 11$.

\begin{table}[H]
    \centering
    \renewcommand{\arraystretch}{1.5}
    \begin{tabular}{|c|c|c|c|c|}
        \hline
        \rowcolor{gray!20}
        \diagbox{$Q_1 Q_0$}{$Q_3 Q_2$} & \textbf{00} & \textbf{01} & \textbf{11} & \textbf{10} \\ \hline
        \textbf{00} & $0$ & $0$ & $0$ & $0$ \\ \hline
        \textbf{01} & $0$ & $0$ & $0$ & $0$ \\ \hline
        \textbf{11} & $0$ & $1$ & $0$ & $0$ \\ \hline
        \textbf{10} & $0$ & $0$ & $0$ & $0$ \\ \hline
    \end{tabular}
    \caption{Harta K pentru ieșirea ERR\_LED}
\end{table}

Ecuația extrasă:
\begin{equation}
    ERR\_LED = \overline{Q_3} Q_2 Q_1 Q_0
\end{equation}

\subsection{Implementarea ieșirilor de recepție}

Aceste semnale activează difuzorul și indicatorul vizual de recepție. Sunt active în starea S11, unde $Q_3 Q_2 = 10$ și $Q_1 Q_0 = 11$.

\begin{table}[H]
    \centering
    \renewcommand{\arraystretch}{1.5}
    \begin{tabular}{|c|c|c|c|c|}
        \hline
        \rowcolor{gray!20}
        \diagbox{$Q_1 Q_0$}{$Q_3 Q_2$} & \textbf{00} & \textbf{01} & \textbf{11} & \textbf{10} \\ \hline
        \textbf{00} & $0$ & $0$ & $0$ & $0$ \\ \hline
        \textbf{01} & $0$ & $0$ & $0$ & $0$ \\ \hline
        \textbf{11} & $0$ & $0$ & $0$ & $1$ \\ \hline
        \textbf{10} & $0$ & $0$ & $0$ & $0$ \\ \hline
    \end{tabular}
    \caption{Harta K pentru SPK\_EN și RX\_LED}
\end{table}

Ecuația extrasă:
\begin{equation}
    SPK\_EN = RX\_LED = Q_3 \overline{Q_2} Q_1 Q_0
\end{equation}

\subsection{Implementarea ieșirilor de transmisie}

Aceste semnale activează microfonul și indicatorul vizual de emisie. Sunt active în starea S13, unde $Q_3 Q_2 = 11$ și $Q_1 Q_0 = 01$.

\begin{table}[H]
    \centering
    \renewcommand{\arraystretch}{1.5}
    \begin{tabular}{|c|c|c|c|c|}
        \hline
        \rowcolor{gray!20}
        \diagbox{$Q_1 Q_0$}{$Q_3 Q_2$} & \textbf{00} & \textbf{01} & \textbf{11} & \textbf{10} \\ \hline
        \textbf{00} & $0$ & $0$ & $0$ & $0$ \\ \hline
        \textbf{01} & $0$ & $0$ & $1$ & $0$ \\ \hline
        \textbf{11} & $0$ & $0$ & $0$ & $0$ \\ \hline
        \textbf{10} & $0$ & $0$ & $0$ & $0$ \\ \hline
    \end{tabular}
    \caption{Harta K pentru MIC\_EN și TX\_LED}
\end{table}

Ecuația extrasă:
\begin{equation}
    MIC\_EN = TX\_LED = Q_3 Q_2 \overline{Q_1} Q_0
\end{equation}

\vspace{1cm}
\begin{tikzpicture}[
    circuit logic US,
    small circuit symbols, % Reduced gate size
    every node/.style={font=\footnotesize},
]

% =========================================================
% 1. Error Output (S7)
% =========================================================
\node[anchor=west] at (-1, 3.75) {\textbf{1. Error Output (S7)}};

% The 4-input AND Gate positioned at x=3
\node[and gate, inputs={nnnn}, point right] (AND1) at (3,1.875) {};

% Inputs positioned at x=0 with generous vertical spacing (0.4cm gap)
\node (q3_1) at (0, 3) {$Q_3$};
\node (q2_1) at (0, 2.25) {$Q_2$};
\node (q1_1) at (0, 1.5) {$Q_1$};
\node (q0_1) at (0, 0.75) {$Q_0$};

% Wires & NOT Gate
% Q3 -> NOT -> Input 1
\node[not gate, point right] (NOT1) at (1.5, 3) {}; % NOT gate in the middle (x=1.5)
\draw (q3_1) -- (NOT1.input);
\draw (NOT1.output) -- ++(right:0.2) |- (AND1.input 1);

% Q2 -> Input 2 (Direct wire, bending into input)
\draw (q2_1) -- ++(right:1.2) |- (AND1.input 2);

% Q1 -> Input 3
\draw (q1_1) -- ++(right:1.2) |- (AND1.input 3);

% Q0 -> Input 4
\draw (q0_1) -- ++(right:1.5) |- (AND1.input 4);

% Output Label
\draw (AND1.output) -- ++(0.5,0) node[right] {\textbf{ERR\_LED}};


% =========================================================
% 2. Reception Output (S11)
% =========================================================
\node[anchor=west] at (-1, -0) {\textbf{2. Reception Output (S11)}};

\node[and gate, inputs={nnnn}, point right] (AND2) at (3,-1.875) {};

% Inputs for Circuit 2
\node (q3_2) at (0, -0.75) {$Q_3$};
\node (q2_2) at (0, -1.5) {$Q_2$};
\node (q1_2) at (0, -2.25) {$Q_1$};
\node (q0_2) at (0, -3) {$Q_0$};

% Q3 -> Input 1
\draw (q3_2) -- ++(right:2.25) |- (AND2.input 1);

% Q2 -> NOT -> Input 2
\node[not gate, point right] (NOT2) at (1.25, -1.5) {};
\draw (q2_2) -- (NOT2.input);
\draw (NOT2.output) -- ++(right:0.2) |- (AND2.input 2);

% Q1 -> Input 3
\draw (q1_2) -- ++(right:1.25) |- (AND2.input 3);

% Q0 -> Input 4
\draw (q0_2) -- ++(right:1.5) |- (AND2.input 4);

% Output Label
\draw (AND2.output) -- ++(0.5,0) node[right, align=left] {\textbf{SPK\_EN} \\ \textbf{RX\_LED}};


% =========================================================
% 3. Transmission Output (S13)
% =========================================================
\node[anchor=west] at (-1, -3.75) {\textbf{3. Transmission Output (S13)}};

\node[and gate, inputs={nnnn}, point right] (AND3) at (3,-5.625) {};

% Inputs for Circuit 3
\node (q3_3) at (0, -4.5) {$Q_3$};
\node (q2_3) at (0, -5.25) {$Q_2$};
\node (q1_3) at (0, -6) {$Q_1$};
\node (q0_3) at (0, -6.75) {$Q_0$};

% Q3 -> Input 1
\draw (q3_3) -- ++(right:2.25) |- (AND3.input 1);

% Q2 -> Input 2
\draw (q2_3) -- ++(right:1.25) |- (AND3.input 2);

% Q1 -> NOT -> Input 3
\node[not gate, point right] (NOT3) at (1.25, -6) {};
\draw (q1_3) -- (NOT3.input);
\draw (NOT3.output) -- ++(right:0.2) |- (AND3.input 3);

% Q0 -> Input 4
\draw (q0_3) -- ++(right:2.25) |- (AND3.input 4);

% Output Label
\draw (AND3.output) -- ++(0.5,0) node[right, align=left] {\textbf{MIC\_EN} \\ \textbf{TX\_LED}};

\end{tikzpicture}
